\subsection{The differentail in the enriched model}
until the end of this seciton, we will develope the tools to describe the differential, which is the the composition in \cref{lem:Ch_diff}, in the simplicall and topological enriched model.

sketch:
define the cube category which is the nerve
\[
    N 1^n = (\Delta^1)^n
\]

define a spine (and mention that there is "the" spine but all of them give the same result).
define spine functors in a category C (modeled as quasi category).
Remind that there is an adjunction Nerve and rigidification.
Take D to be sSet enriched model for C, i.e. N D = C.
we get:
\[
    \Fun^{\Spine} (cube_n, N D) = \Fun^{\Spine} (\CC cube_n, D)
\]

So spine functors from $\CC cube_n$ are the sSet-Cat model for chain complexes.

Note that we acn also define the pointed category, by adding a object $*$ to $cube_n$ with unique morphisms
\[
    * \to x, \quad x \to *
\]
and for any $x \notin \Spine_n$ we invert (localize) the map $x \to *$.
This give equivalence:
\[
    cube_n / \{x \to * \mid x \notin \Spine_n\} \simeq \Ch_n
\]
which given by
\[
    x \mapsto \begin{cases*}
        *               & $x \notin \Spine_n$ \\
        \lvert x \rvert & $x \in \Spine_n$
    \end{cases*}
\]

Given two maps of n chain complexes in $\cC$, the form:
\[
    X_n \xrightarrow{f} K \xrightarrow{g} Y_0
\]
where the functors (chain complexes) $X_n, Y_n$ are zero except on $n$ and $0$, we can use the sSet model for functor categories to get two cubes in $\cC$
\[
    f,g \in \Fun(\Delta \times N 1^n, \cC)
\]
such that:
\[
    \partial_0 f = \partial_1 g = K.
\]
where $\partial_i$ is the restriction to $\partial_i \Delta^1$ in the first coordinate.
We can glue to get the following horn in $\cC$
\[
    f\cup g : \Lambda^2_1 \times N 1^n \to \cC
\]

Having that $\cC$ is quasi-category, we get a horn filling:
\[
    \Delta^2 \times N 1^n \to \cC
\]

we denote the restriction
\[
    g \circ f : \partial_1 \Delta^2 \times N 1^n \to \cC
\]

Indeed, as a morphism:
\[
    g \circ f : X_n \to Y_0
\]
hence as a n+1 cube in $\cC$ we have:
\[
    g \circ f (A) = \begin{cases*}
        X & $A = n+1$ \\
        Y & $A = 0$   \\
        0 & otherwise
    \end{cases*}
\]

In the Cat-sSet model we get that $g \circ f$ is geiven by a map of Cat-sSet
\[
    \CC (\partial_1 \Delta^2 \times N 1^n) = \CC (N 1^{n+1}) \to D
\]
All the object except the initial and the terminal are maps to $0 \in D$ trivail, hence the only non tirival mapping sSet is:
\[
    \Hom_{\CC(\hN 1^{n+1})}(1_{n+1}, 0_{n+1}) \to \Hom(X,Y)
\]

\begin{defn}[Simplicial Mapping Spaces]
    For a general poset $A$ and elements $a \le b$, the simplicial mapping space in the rigidified nerve $\mathfrak{C}[N(A)]$ is the nerve of the poset of strict inclusion chains:
    \[ \operatorname{Hom}_{\mathfrak{C}[N(A)]}(a, b) = N(\mathcal{P}_{a,b}) \]
    where $\mathcal{P}_{a,b}$ consists of chains $a = X_0 \subsetneq X_1 \subsetneq \dots \subsetneq X_m = b$ \todo{ADD REF??}. The higher $k$-simplices of this mapping space are given by sequences of chain inclusions $C_0 \subseteq C_1 \subseteq \dots \subseteq C_k$ \todo{ADD REF??}.
\end{defn}

\begin{defn}[The Permutohedron]
    Let $A$ be a finite linearly ordered set, and let $Cu_A$ be the cube category modeled as the Boolean lattice of subsets of $A$. The mapping space from the initial object $\emptyset$ to the maximal object $A$ is called the simplicial permutohedron $P_A$ \todo{ADD REF??}:
    \[ P_A := \operatorname{Hom}_{\mathfrak{C}[N(Cu_A)]}(\emptyset, A) = N(\mathcal{P}_{\emptyset, A}) \]
\end{defn}

\begin{defn}[Composition and Boundary]
    For any proper intermediate subset $Q$ ($\emptyset \subsetneq Q \subsetneq A$), the composition of morphisms through $Q$ is defined by the union of chains \todo{ADD REF??}:
    \[ \circ : \mathcal{P}_{Q,A} \times \mathcal{P}_{\emptyset,Q} \to \mathcal{P}_{\emptyset,A}, \quad (C_1, C_2) \mapsto C_1 \cup C_2 \]
    We define the geometric boundary $\partial P_A$ (a subcomplex of dimension $\dim(P_A) - 1$) as the union over all proper intermediate subsets $Q$ of the images of these composition embeddings \todo{ADD REF??}:
    \[ \partial P_A := \bigcup_{\emptyset \subsetneq Q \subsetneq A} P_{A \setminus Q} \times P_Q \]
\end{defn}

\begin{fact}[Cone Structure and Contractibility]
    The poset of chains $\mathcal{P}_{\emptyset, A}$ contains a unique minimal element: the empty chain $C_{min} = (\emptyset \subsetneq A)$, which corresponds to the direct, unfactored jump from source to target \todo{ADD REF??}.

    Because $\mathcal{P}_{\emptyset, A}$ has a unique minimum, its nerve $P_A$ is topologically a cone over the sub-poset of proper intermediate chains $P^\circ$ \todo{ADD REF??}. In simplicial terms, $P_A$ is the join of the 0-simplex (the cone point) and its boundary:
    \[ P_A \cong Cone(N(P^\circ)) \cong \Delta^0 \star N(P^\circ) \]
    Because $P_A$ is formed as a cone, it is trivially contractible \todo{ADD REF??}.
\end{fact}


To make things clean, we reindex the source of $f,g$ according to the horn:
\[
    \Lambda^2_1 \times N 1^n \to \cC
\]
we name
\begin{gather*}
    0 \coloneq (0,\varnothing) \\
    1 \coloneq (1,\varnothing) \\
    n+1 \coloneq (1, \n) \\
    n+2 \coloneq (2,\n) 
\end{gather*}
note that under this "change of variables" we have that $K$ is located at ${1} \times N 1^n$ and  $1, n+1$ are the first and the last elements of $K$. By abuse of notation we will write:
\[
    K_{n+2} = X, \quad K_0 = Y
\]

Let $n+2 > i > 1$ we denote
\[
     f_i : \Hom_{\CC N 1^{n+1}}(n+2,i) \to \Hom(X,K_i)
\]
and for $n+1 > j > 0$ we denote
\[
    g_j: \Hom_{\CC N 1^{n+1}}(j,0) \to \Hom(K_j,Y)
\]
to be the maps that induced by modeling the functors $f,g$ in the Cat-sSet model.
Similarly, by using the Cat-sSet model, for the functor $K$ we get maps
\[
    K_{i,j} : \Hom_{\CC N 1^n}(i,j) \to \Hom(K_i,K_j)
\]